\chapter{1956:Sir Cyril Hinshelwood and Nikolay Semenov}

\label{chap:chemistry1956}

The Nobel Prize in Chemistry for the year 1956 was awarded jointly to Sir Cyril Hinshelwood and Nikolay Semenov.

\textbf{Motivation:}

for their researches into the mechanism of chemical reactions

\section{Sir Cyril Hinshelwood}

\subsection*{Early Life and Education}

Sir Cyril Hinshelwood was born on 1897-06-19 in London, England.

\subsection*{Career and Major Research}

Hinshelwood studied at Balliol College, Oxford, and spent his career at Oxford, where he became Dr Lee's Professor of Chemistry. He studied the kinetics of gas-phase reactions and showed that many of them proceed by chain mechanisms. His investigations of the explosive reaction between hydrogen and oxygen established the concepts of chain branching and explosion limits, and he went on to study the kinetics of bacterial growth.

\subsection*{Nobel Prize-winning Work}

The work for which Sir Cyril Hinshelwood was awarded the Nobel Prize in Chemistry in 1956 was described by the Nobel Committee as: "for their researches into the mechanism of chemical reactions".

Hinshelwood showed how the rates of chemical reactions could be interpreted in terms of the sequence of elementary steps, and he demonstrated that chain reactions with branched propagation could account for the sharp explosion limits of the hydrogen-oxygen reaction.

\subsection*{Significance and Impact}

Hinshelwood's kinetic analysis provided a quantitative framework for understanding reaction mechanisms, especially chain reactions, and connected the theory of gas reactions to practical phenomena such as explosions. His later work extended kinetic methods to the growth of bacteria, linking chemistry to biology.

\subsection*{Later Career and Honors}

Hinshelwood continued at Oxford and later served as president of the Royal Society of London. He was knighted in 1948 and created an officer of the Order of Merit. He died on 1967-10-09 in London.

\subsection*{Legacy}

The chain-reaction theory developed by Hinshelwood and Semenov became a cornerstone of chemical kinetics, and Hinshelwood's name remains associated with the quantitative study of reaction mechanisms in chemistry.

\subsection*{Selected Publications}

\begin{itemize}

\item \emph{The Kinetics of Chemical Change in Gaseous Systems}, Oxford: Clarendon Press, 1926.

\end{itemize}

\clearpage

\section{Nikolay Semenov}

\subsection*{Early Life and Education}

Nikolay Semenov was born on 1896-04-15 in Saratov, Russia.

\subsection*{Career and Major Research}

Semenov studied physics and chemistry at Petrograd University and worked in Leningrad, where he founded and directed the Institute of Chemical Physics of the Soviet Academy of Sciences, later moved to Moscow. Independently of Hinshelwood, he developed the theory of chain reactions and in 1934 published his book \emph{Chemical Kinetics and Chain Reactions}. He showed that branched chain reactions could account for the ignition and explosion limits observed in the oxidation of phosphorus and hydrogen.

\subsection*{Nobel Prize-winning Work}

The work for which Nikolay Semenov was awarded the Nobel Prize in Chemistry in 1956 was described by the Nobel Committee as: "for their researches into the mechanism of chemical reactions".

Semenov's theory of branched chain reactions explained how a single reactive intermediate could multiply and accelerate a reaction until it becomes explosive, and it provided the mathematical treatment of the explosion limits of the hydrogen-oxygen and phosphorus-oxygen reactions.

\subsection*{Significance and Impact}

Semenov's work established the quantitative theory of chain reactions and combustion, with direct applications to flames, explosions, and industrial chemical processes. He became a central figure in Soviet physical chemistry, training a large school of kineticists.

\subsection*{Later Career and Honors}

Semenov directed the Institute of Chemical Physics in Moscow for the rest of his career and played a leading role in the organization of science in the Soviet Union. He died on 1986-09-25 in Moscow.

\subsection*{Legacy}

Semenov is regarded as one of the founders of chemical kinetics, and his theory of branched chain reactions is fundamental to the modern understanding of combustion and explosions.

\subsection*{Selected Publications}

\begin{itemize}

\item \emph{Chemical Kinetics and Chain Reactions}, Oxford: Clarendon Press, 1935.

\end{itemize}

\clearpage

\section*{Chapter Summary}

The 1956 prize recognized research into the mechanism of chemical reactions. Working independently, Sir Cyril Hinshelwood at Oxford and Nikolay Semenov in the Soviet Union developed the theory of chain reactions and showed how branched chain mechanisms explain the explosion limits of gas-phase reactions such as the oxidation of hydrogen.

